% Proposed response to Comment 3, using the agreed saved-baseline figures.
% Upload decay_constant_ratio_vs_theta.pdf into the reply project.
% The reply preamble needs graphicx and amsmath.
% Correct the manuscript Figure 1/Figure 2 image swap before submission.
% This fragment does not claim that the Figure 3 caption has been edited.
% No OPE correction or numerical baseline replacement is implied.

\noindent\textbf{Response:}

We thank the referee for highlighting the independent interest of the
decay constants and their potential application to nonleptonic $B$ decays.
We have added Fig.~1 in Sec.~III of the revised manuscript, showing
$f_1(\theta)$ and $f_2(\theta)$ separately for both the charm-strange
and bottom-strange sectors. This displays the individual quantities
entering our radiative-decay sum rules. The corresponding ratio
$f_2(\theta)/f_1(\theta)$ requested by the referee is presented in
Fig.~\ref{fig:referee1-ratio} of this response.

Both figures use the same representative QCD inputs, two-point Borel
parameters, and continuum thresholds as the angular scan of the radiative
widths. The decay constants are recalculated at each angle, and the ratio
is formed from their values at that same angle. For these fixed input
choices, we obtain
\begin{equation}
\left.\frac{f_2^{(D_s)}}{f_1^{(D_s)}}\right|_{\theta=26.6^\circ}
\simeq0.418,
\qquad
\left.\frac{f_2^{(B_s)}}{f_1^{(B_s)}}\right|_{\theta=38.5^\circ}
\simeq0.166.
\end{equation}
These are illustrative values at fixed inputs, rather than new estimates
with propagated uncertainties. The radiative-width angular scan also
includes the angle dependence of both decay constants.

For the suggested factorization application, we clarify that the $f_i$
in our manuscript are defined through the rotated interpolating currents.
The weak axial-current constants in the charm-strange sector, defined by
$\langle0|\bar s\gamma_\mu\gamma_5 c|A_i\rangle
=m_{A_i}f_{A,i}\eta_\mu$, are related to these quantities within the
adopted assumption that each rotated current couples to its corresponding
physical state. The inverse current rotation gives
\begin{equation}
f_{A,1}=\sin\theta\,f_1^{(D_s)},\qquad
f_{A,2}=\cos\theta\,f_2^{(D_s)},\qquad
\frac{f_{A,2}}{f_{A,1}}=\cot\theta\,
\frac{f_2^{(D_s)}}{f_1^{(D_s)}}.
\end{equation}
The nonleptonic branching-fraction ratio additionally requires the
$B\to D$ transition form factors and phase-space factors. A numerical
analysis of those decays is beyond the scope of the present radiative
study.

% This numbering applies to figures in the referee-response document only.
\renewcommand{\thefigure}{R\arabic{figure}}
\begin{figure}[tbp]
  \centering
  \includegraphics[width=\linewidth]{decay_constant_ratio_vs_theta.pdf}
  \caption{Mixing-angle dependence of the ratio $f_2(\theta)/f_1(\theta)$
  of the rotated-current decay constants for the (a) charm-strange and
  (b) bottom-strange sectors. Both constants are recalculated at each
  angle with the same fixed representative inputs used for Fig.~1
  of the revised manuscript. Red dash-dotted lines mark the adopted
  central angles $26.6^\circ$ and $38.5^\circ$ in the respective panels;
  gray dotted lines mark the heavy-quark-limit reference $35.3^\circ$.
  The curves illustrate the angular dependence at fixed inputs;
  no uncertainty bands are shown.}
  \label{fig:referee1-ratio}
\end{figure}
